\section{Lifted shifted correlations}\label{sec:shifts}

We retain the prime groups and their notation from
\Cref{sec:transference,sec:ideal}. Thus the small groups have primes between
$\exp(L^a)$ and $\exp(L^b)$, the big groups are intervals of primes between
$\exp(L^c)$ and $\exp(L^d)$, and
$0<a<b<c<d<0.47$. All groups are disjoint, their harmonic masses
$V_g$ lie in $[v_-,v_+]$, and each kind has between fixed positive
multiples of $T$ groups. The parameters $q\in(0,1)$ and $\ell\geq1$
are fixed. Write $\mathcal P=\bigcup_g\mathcal P_g$ and let $n_*$ be the
part of $n$ supported on primes outside $\mathcal P$.

\subsection{Endpoint hypotheses and the correlation theorem}

For a vector $\mathbf t=(t_g)_g$ of nonnegative integers, define
\begin{equation}\label{eq:marked-weight}
 W_{\mathbf t}(n)=q^{\omega(n)-\sum_g t_g}
       \prod_g\frac{(\omega_g(n))_{t_g}}{V_g^{t_g}}.
\end{equation}
The falling factorial counts ordered lists of distinct group primes
dividing $n$. More generally, if $\mathbf b$ is a function on these lists,
put
\[
 W_{\mathbf t}^{\mathbf b}(n)=
 q^{\omega(n)-\sum_g t_g}
 \prod_g V_g^{-t_g}
 \sum_{\substack{(p_{g,1},\ldots,p_{g,t_g})\text{ distinct in each group}\\
                  p_{g,j}\in\mathcal P_g,\ p_{g,j}\mid n}}
                  \mathbf b((p_{g,j})_{g,j}).
\]
An empty admissible-list set gives value zero. Write $W_\ell$ when all
$t_g=\ell$. For each fixed $t_{\max}$,
\begin{equation}\label{eq:mark-sup}
 \sup_{\max t_g\leq t_{\max}}
       |W_{\mathbf t}^{\mathbf b}(n)|
 \leq \|\mathbf b\|_\infty L^{C(t_{\max})}.
\end{equation}
Indeed, $q^{y-t}(y)_t/V_g^t$ has a bounded supremum over integers $y\geq t$,
uniformly for $t\leq t_{\max}$, and there are $O(T)$ groups.

An endpoint is a function invariant under multiplication by
$\mathcal P$-integers of the form
\begin{equation}\label{eq:endpoint-form}
 F_0(n)=\sum_{mpr=n_*}
   \alpha_m\chi_0(m)m^{i\sigma_0}
   \ind_{p\in I_p,\,\Pminus(p)>W}\,p^{i\sigma_1}c_r.
\end{equation}
Here $|\alpha_m|,|c_r|\leq L^C$, $\alpha_m=0$ unless
$\Pminus(m)>W$, $I_p\subset[x^\tau,x^\eta]$ is an interval,
$0<\tau<\eta<1/4$ are fixed, $\chi_0$ is a Dirichlet character of
modulus at most $L^C$, and $|\sigma_0|,|\sigma_1|\leq L^C$.
The variable $p$ ranges over integers. The sequences may depend on $x$,
and all the parameters and sequences can differ at the two endpoints.
Since $W=\exp(\sqrt L)$ exceeds every
group prime, $m$ and $p$ are $\mathcal P$-free for sufficiently large $x$.
We impose the following condition on the $\alpha$ of at least the first
endpoint:
\begin{equation}\label{eq:discrepancy}
 \left|\sum_{m\in I}\frac{\alpha_m\chi(m)m^{iu}}m\right|
       \leq L^{-B}
 \quad\left\{
 \begin{array}{l}
 I\subset[1,x^2]\text{ an interval},\\
 \chi\text{ a Dirichlet character of modulus at most }L^B,\\
 |u|\leq L^B.
 \end{array}\right.
\end{equation}
This includes principal and imprimitive characters. An existing range
restriction on $\alpha$ is part of the sequence in this condition.

\begin{theorem}[Lifted shift cancellation]\label{thm:shift}
Fix all the group and endpoint bounds above. Fix $D_0,C>0$ and a fixed
compact interval $[c_\Psi,C_\Psi]\subset(0,\infty)$. Suppose
$0<|k|\leq L^C$ is integral, $xL^{-C}\leq Z\leq xL^C$,
$|a_1|,|a_2|\leq L^C$, and $\Psi$ is supported in that interval with
$\|\Psi^{(j)}\|_\infty\ll_j L^{C(j+1)}$ for every $j\geq0$.
There exists a fixed $B$, depending only on these bounds and $D_0$,
such that \Cref{eq:discrepancy} implies
\begin{equation}\label{eq:shift-bound}
 \sum_{\substack{z>0\\z+k>0}}
 \frac{\Psi(z/Z)}z\,z^{ia_1}(z+k)^{ia_2}
       \overline{F_0(z)}G_0(z+k)W_\ell(z)W_\ell(z+k)
       \ll L^{-D_0}.
\end{equation}
The bound is uniform over the indicated sequences, intervals, characters,
frequencies, and cutoffs. It also holds when the discrepancy-bearing
endpoint or its coefficients have been conjugated.
\end{theorem}

Conjugation in the last assertion is harmless: conjugate
\Cref{eq:discrepancy} and replace $(\chi,u)$ by $(\overline\chi,-u)$.
Products with additional characters and bounded power twists are covered
by increasing $B$. We give the proof in several stages.

\subsection{The physical lift and removal of shared labels}

We first show that the raw correlation together with its signed
comparison correlations is small. The harmonic factor $1/n$ will
turn sums over shared prime divisors into draws with laws $\mu_g$
when the shared product is removed.

Take the raw pattern and signed comparison patterns of \Cref{thm:ideal}.
Write $m_{\rm probe}$ for the fixed probe bound $m$ in that theorem;
the letter $m$ in \Cref{eq:endpoint-form} denotes an integer factor.
As in the physical construction, $M=J+\ell$, the edge product is $D$,
and the positions are $n,n'=n+kD$. Insert into the physical pairing the
root weight $\Psi(n/(ZD))/n$, the edge phase $D^{-i(a_1+a_2)}$, and
the mark-independent vectors
\begin{equation}\label{eq:lift-vectors}
 f(n,\mathbf p)=n^{-ia_1}F_0(n)q^{(\omega(n)-Ms)/2},
 \qquad
 g(n',\mathbf p')=(n')^{ia_2}G_0(n')q^{(\omega(n')-Ms)/2}.
\end{equation}
Our pairing is conjugate-linear in its first entry, so its integrand
contains $n^{ia_1}\overline{F_0(n)}$ at the source and
$(n')^{ia_2}G_0(n')$ at the target.
Every factor involving $D$ is inside the slot symmetrizations.
On the support, $n\asymp ZD=x^{1+o(1)}$ and
$n'/n=1+O(L^C/Z)$. A partition into at most $O(L)$ dyads
$n\asymp X$ therefore suffices, and one can restrict the other endpoint
to an interval of length $O(X)$.

We check carefully that the exponent in the endpoint norm bound is
independent of $J$ and of the probe count. On a physical state
$\omega(n)\geq Ms$, so the damping in \Cref{eq:lift-vectors} is at most
one. There are at most $\tau(n_*)^2$ ordered factorizations in
\Cref{eq:endpoint-form}, whence
$|F_0(n)|\leq L^{2C}\tau(n_*)^2$. Fix an ordered physical list
$S'$ of $M$ distinct labels per group, with product $P(S')$.
For $P(S')\mid n$, one has $n_*\mid n/P(S')$ and hence
\[
 \sum_{\substack{n\asymp X\\P(S')\mid n}}\tau(n_*)^4
 \leq\sum_{v\asymp X/P(S')}\tau(v)^4
 \ll\frac{X}{P(S')}L^{C_5}.
\]
Here $P(S')\leq\exp(O_J(TL^d))=x^{o(1)}$, so the interval in $v$
is still of positive power length. Moreover
\[
 \sum_{S'}\frac1{P(S')}\prod_g V_g^{-M}
 \leq\prod_g\left(\sum_{p\in\mathcal P_g}\frac1{pV_g}\right)^M=1.
\]
Thus both endpoint vectors have norm at most $X^{1/2}L^{C_6}$ in the
physical state measure, with $C_6$ independent of $J$ and the probe count.
Also $|F_0(n)|\leq\exp(O(\sqrt L))$: an integer of size $x^{O(1)}$
has $O(\sqrt L)$ prime factors above $W$, with multiplicity, so it has
at most $2^{O(\sqrt L)}$ choices for each of the rough divisors $m,p$.
The vectors are independent of the selected marks.

For clarity, let $\Phi(u)=\Psi(e^u)$. Logarithmic Fourier inversion gives
\[
 \Psi(n/(ZD))=\frac1{2\pi}\int_{\R}
       \widehat\Phi(v)(n/Z)^{iv}D^{-iv}\,dv.
\]
The derivative hypotheses imply an $L^{O(1)}$ integral norm, and, for a
fixed sufficiently large $C_7$, tails $|v|>L^{C_7}$ have arbitrarily
large negative log powers. This cutoff exponent can be fixed from the
derivative bounds: increasing the number of integrations by parts
increases the saving without increasing $C_7$.
On the dyad multiply the first vector by $X/n$ and $(n/Z)^{-iv}$,
and divide the pairing by $X$. The conjugation in the pairing then
supplies the factor $(n/Z)^{iv}$ above. The required ideal frequency is
$\zeta=-a_1-a_2-v$. Consequently \Cref{cor:transference-pairing}
and \Cref{thm:ideal} bound the entire residual pairing by $L^{-D_1}$,
for any prescribed $D_1$, after choosing the transfer target first,
then the ideal family, then $J$. The norm exponent just proved makes
this ordering possible. The discarded Fourier tails obey the same
conclusion by \Cref{lem:physical-schur}; their accuracy may be chosen
after the absolute comparison costs are known. All these steps apply
to the sum of the patterns, with their signs.

We next calculate the individual pattern after this lift. In group $g$
write $c_g$ for its free-slot count and $s_{R,g}=J-c_g$ for its
shared-slot count; thus $t_g=\ell+c_g$. Let $D_R,D_C$ be the products
of the shared and free labels, so $D=D_RD_C$. This partitions labels
by their role on the edge; the earlier factorization
$D=D_{\mathcal S}D_{\mathcal B}$ partitions them by prime-group size.
Set
$n=D_Rw$, $n'=D_Rw'$, so $w'=w+kD_C$.
The physical state and target-list normalizations in this group are
\[
 V_g^{-M-t_g}=V_g^{-s_{R,g}-2t_g},
 \qquad M+t_g=s_{R,g}+2t_g.
\]
The factor $1/D_R$ from $1/n$ therefore changes each shared-label
sum into a draw of mass $1/(pV_g)$, and leaves $V_g^{-t_g}$ at each
endpoint for its unshared lists. Away from shared overlaps,
\[
 \omega_g(D_Rw)-M=\omega_g(w)-t_g.
\]
The damping already in the graph and that in \Cref{eq:lift-vectors}
give this full power of $q$ at each endpoint. Invariance gives
$F_0(D_Rw)=F_0(w)$, and the phases become
\[
 (D_Rw)^{ia_1}(D_Rw')^{ia_2}(D_RD_C)^{-i(a_1+a_2)}
       =(w/D_C)^{ia_1}(w'/D_C)^{ia_2}.
\]
The coefficient $K_\nu$ of a comparison uses the free labels and
designated unshared big labels, and uses no shared labels. Thus its
dependence is preserved when these shared harmonic draws are summed.

Here are error bounds justifying the exclusions in this calculation.
For fixed $w,w',D_C$ and any earlier shared draws, only $O_J(L)$
values are forbidden to a new shared draw: endpoint prime divisors,
free labels, and previous shared labels. Each atom has mass at most
$v_-^{-1}\exp(-L^a)$. The $O_J(T)$ draws have total excluded
probability $L^{O_J(1)}\exp(-L^a)$. For distinct shared values the exact
identity is
\[
 \omega(D_Rw)-Ms=\omega(w)-\sum_g t_g
             -\#\{p:p\mid D_R,\ p\mid w\}.
\]
In bounding the actual overlap terms by the unrestricted marked sum,
the damping loss is at most $q^{-2Ms}=L^{O_J(1)}$. Unshared divisors
still divide $w,w'$ and their count cannot increase. By
\Cref{eq:mark-sup}, endpoint Cauchy--Schwarz and divisor moments, their
remaining absolute harmonic average is $L^{O(1)}$. Shared exclusions
therefore cost $L^{O_J(1)}\exp(-L^a)$.

Once shared overlaps are excluded, the remaining ban excludes free
labels from the unshared endpoint mark lists. Every excluded
configuration therefore has a free prime $p'$ dividing an endpoint.
Since $p'\mid D_C$, it follows that both $p'\mid w$ and $p'\mid w'$;
it suffices to bound all such multiples. On $w\asymp Y_0=ZD_C$, put $w=p'v$ and
$w'=p'(v+kD_C/p')$. Invariance and Cauchy--Schwarz give
\[
 \sum_{\substack{w\asymp Y_0\\p'\mid w}}
 \frac{|F_0(w)G_0(w+kD_C)|}{w}
 \ll \frac{L^{O(1)}}{Y_0}
 \left(\sum_{v\asymp Y_0/p'}\tau(v)^4\right)^{1/2}
 \left(\sum_{v\asymp Y_0/p'}
            \tau(v+kD_C/p')^4\right)^{1/2}
 \ll\frac{L^{O(1)}}{p'}.
\]
Both shifted ranges stay in positive intervals of size $O(Y_0/p')$,
and $Y_0/p'=x^{1+o(1)}$. The mark weights, coefficients, and total
pattern costs are fixed log powers. Summing over the free slots
therefore makes this error negligible too. This is an averaged bound
over multiples of $p'$, not a restriction on the arbitrary coefficients.

We have obtained, up to errors smaller than every fixed negative log
power, the following expression for each pattern:
\begin{equation}\label{eq:stripped-pattern}
 \E_{\rm free}\sum_{\substack{w>0\\w'=w+kD_C>0}}
 \frac{\Psi(w/(ZD_C))}{w}
 (w/D_C)^{ia_1}(w'/D_C)^{ia_2}
 \overline{F_0(w)}G_0(w')
 W_{\mathbf t}(w)W_{\mathbf t}(w')\,
 \E_{\rm marks}K_\nu.
\end{equation}
The outer expectation uses independent free-label laws $\mu_g$.
Conditional on these labels and the positions, the inner expectation
uses independent uniform ordered unshared lists at each endpoint;
its term is zero if either list set is empty. The raw pattern has
$D_C=1$ and $t_g=\ell$, so it is exactly the sum in
\Cref{eq:shift-bound}. It remains to bound all the comparison terms.

\subsection{Comparison shifts and their major arcs}

Expand a comparison coefficient from \Cref{lem:comparison-kernel} by
its two logarithmic cells and two characters. Its factors on the free
products $D_A,D_B$ separate, and the remaining factors are bounded
weights on specified big unshared marks at the two endpoints. Cell
normalizations, the number of terms, and the integral costs are fixed
log powers, independent of $J$. Let $H'$ be the product of the lower
size endpoints of the two cells. Then $D_AD_B\asymp H'$ and
\begin{equation}\label{eq:comparison-scales}
 H'\geq\exp(L^c),\qquad
 H'\leq\exp(O(m_{\rm probe}TL^d)),\qquad
 Y=ZH'=x^{1+o(1)}.
\end{equation}
Insert smooth size cutoffs at $w,w'\asymp Y$, separate
$\Psi(w/(ZD_AD_B))$ logarithmically, and include $Y/w$ in the first
cutoff. Apart from $1/Y$ and fixed log-power costs, every resulting
term is
\begin{equation}\label{eq:comparison-correlation}
 \E_{D_A,D_B}b_1(D_A)b_2(D_B)
         \sum_w\overline{\mathcal U(w)}\mathcal V(w+kD_AD_B).
\end{equation}
Here $|b_1|,|b_2|\leq1$ after extracting constants; each is supported
on its cell. Each endpoint sequence is its invariant endpoint times a
smooth cutoff at size $Y$, a power $w^{iv}$ with $|v|\leq L^{O(1)}$,
and a weight $W_{\mathbf t}^{\mathbf b}(w)$. The additional mark
function uses only big labels. In particular
\begin{equation}\label{eq:comparison-l2}
 \sum_w|\mathcal U(w)|^2+\sum_w|\mathcal V(w)|^2
        \ll YL^{C_8}.
\end{equation}
All frequency and cutoff bounds here are fixed before $B$.

The endpoint sequences are now independent of the free products.
Fourier inversion isolates the weighted shift average as a multiplier:
bilinear cancellation
will control it away from small rational frequencies, and endpoint
Fourier energy will control the remaining contribution.
Use $\widehat F(\theta)=\sum_w F(w)e(\theta w)$ on $\R/\Z$.
The multiplier in \Cref{eq:comparison-correlation} is
\[
 M(\theta)=\E b_1(D_A)b_2(D_B)e(-k\theta D_AD_B).
\]
We claim that, for any fixed desired saving $D_2$, it is
$O(L^{-D_2})$ outside
\[
 \mathfrak M=\bigcup_{1\leq r\leq L^{C_{\rm arc}}}
       \bigcup_{u\bmod r}
       \{\theta:\ |\theta-u/r|_{\R/\Z}
                                  \leq L^{C_{\rm arc}}/H'\},
\]
if $C_{\rm arc}$ is large enough. To verify the coefficient hypothesis
of \Cref{lem:bilinear}, collapse a free product of size $U_0$ to
integer weights $\rho(n)$. Each group contributes at most
$m_{\rm probe}$ slots. Unique factorization bounds the multiplicity
at $n$ by $\prod_g(m_{\rm probe}!)$, and the $V_g$ normalizations
cost another fixed constant per slot. Hence
$\rho(n)\leq L^{O(1)}/n$ and $\sum_n\rho(n)=1$; for bounded tests,
\[
 \sum_{n\asymp U_0}|b(n)\rho(n)|^2\ll L^{O(1)}/U_0.
\]
Both free products are nonempty and each contains a big prime, so
both sizes exceed every fixed power of $L$.

Put $Q_D=\lfloor H'/L^{C'}\rfloor$. Dirichlet approximation to
$k\theta$ gives a reduced $u'/r'$ with $r'\leq Q_D$ and
$|k\theta-u'/r'|\leq1/(r'Q_D)\leq (r')^{-2}$, where
$1/(r'Q_D)\ll L^{C'}/(r'H')$.
If $r'\leq L^{C'}$, division by $k$ places $\theta$ on one of the
stated arcs after increasing $C_{\rm arc}$; reduction can only decrease
the denominator. Otherwise $L^{C'}<r'\ll H'/L^{C'}$.
The normalized bilinear bound of \Cref{lem:bilinear} then saves any
prescribed log power by increasing $C'$. This proves the claim.
Parseval and \Cref{eq:comparison-l2} control the complementary
contribution to \Cref{eq:comparison-correlation}, including its $1/Y$.

Put $H=H'/L^{C_{\rm arc}}$. On each arc it is enough to prove,
for any fixed $A>0$,
\begin{equation}\label{eq:local-energy-target}
 \int_{|\beta|\leq1/H}
       |\widehat{\mathcal U}(u/r+\beta)|^2\,d\beta
          \ll_A YL^{-A}.
\end{equation}
Indeed $|M|\leq1$, and Cauchy--Schwarz with
\Cref{eq:comparison-l2} gives $L^{-A/2+C_8/2}$ per arc after division
by $Y$. There are at most $L^{2C_{\rm arc}+O(1)}$ arcs. We establish
\Cref{eq:local-energy-target} with $A$ chosen after all these costs.
Its proof uses three features of the first endpoint: one small prime
can be extracted from its ordered marks, every factorization contains
the long rough-integer factor in \Cref{eq:endpoint-form}, and the
$m$ coefficients satisfy \Cref{eq:discrepancy}. After conversion to
Mellin frequencies, these supply cancellation in different ranges.

\subsection{One small mark and the rational character expansion}

Choose any small group $g_s$ and a designated slot among its
$t_{g_s}=\ell\geq1$ ordered marks. The endpoint mark weight is
independent of the small labels. Whenever no prime of $\mathcal P_{g_s}$
has its square dividing $w$, removal of the prime in this slot gives
the exact identity
\begin{equation}\label{eq:small-mark-removal}
 W_{\mathbf t}^{\mathbf b}(w)=
 \sum_{\substack{p_s\in\mathcal P_{g_s}\\p_s\mid w}}
 \frac1{V_{g_s}}W_{\mathbf t-\mathbf e_{g_s}}^{\mathbf b}(w/p_s).
\end{equation}
Removing the designated slot is a bijection of the ordered lists;
both $\omega_{g_s}$ and $t_{g_s}$ decrease by one. There is therefore
no extra factorial or damping factor. Both sides, even on the
exceptional integers, are $L^{O(1)}$ by \Cref{eq:mark-sup} and
$\omega_{g_s}(w)=O(L)$. Since
$w_*\mid w/p_s^2$ on a square multiple, the squared norm of the
error after multiplication by $F_0$ is at most
\[
 L^{O(1)}\sum_{p_s\in\mathcal P_{g_s}}
       \sum_{\substack{w\asymp Y\\p_s^2\mid w}}\tau(w_*)^4
 \ll YL^{O(1)}\sum_{p_s\in\mathcal P_{g_s}}p_s^{-2}
 \ll YL^{O(1)}e^{-L^a}.
\]
Every $Y/p_s^2$ is $x^{1+o(1)}$. Parseval makes this error negligible
in \Cref{eq:local-energy-target}.

After \Cref{eq:small-mark-removal}, combine the remaining
$\mathcal P$-part with the residual factor of
\Cref{eq:endpoint-form}. The full integer factorization is
$w=mp_sp r$, with residual coefficient
\[
 c_{r_*}W_{\mathbf t-\mathbf e_{g_s}}^{\mathbf b}(r).
\]
This identity uses that $m,p$ have no group factors. The residual
coefficient is a function of $r$ alone and is bounded by $L^{O(1)}$.

All of $m,p_s,p$ are units modulo the arc denominator $r_0\leq
L^{C_{\rm arc}}$. Separate residual factors by
$d_0=(r,r_0)$ and put $q_0=r_0/d_0$. On units modulo $q_0$ the exact
finite character expansion is
\[
 e(uv/q_0)=\sum_{\chi\bmod q_0}c_{u,q_0}(\chi)\chi(v),
 \qquad
 c_{u,q_0}(\chi)=\frac1{\varphi(q_0)}
           \sum_{v\bmod q_0}^{*}e(uv/q_0)\overline{\chi(v)}.
\]
Each coefficient has modulus at most one. Use this with
$v=mp_sp(r/d_0)$; the condition $(r,r_0)=d_0$ guarantees that $v$
is a unit modulo $q_0$. The factors $\chi(r/d_0)$ and the gcd
restriction are absorbed into the residual coefficient. This
includes every character modulo the actual quotient, so principal
and imprimitive components are retained. The number of divisors and
characters is a fixed log power, and products with $\chi_0$ have
polylogarithmic modulus.

It now suffices to prove the local energy bound at zero for sequences
\begin{equation}\label{eq:four-factor-endpoint}
 a_w=\Psi_1(w/Y)w^{iv_1}
       \sum_{mp_sp r=w}
       \alpha_m\chi_m(m)m^{i\sigma_0}\,
       b_s(p_s)\,
       \ind_{p\in I_p,\,\Pminus(p)>W}\chi_p(p)p^{i\sigma_1}d_r,
\end{equation}
where $p_s\in\mathcal P_{g_s}$, $|b_s|\ll1$, $|d_r|\leq L^{O(1)}$,
the characters and frequencies have fixed log-power bounds, and
$\Psi_1$ has the same kind of smooth size support as above.
Fixed-order divisor moments give $\sum_w|a_w|^2\ll YL^{O(1)}$.

\subsection{From local Fourier energy to a Mellin integral}

We record the scale conversion explicitly. Choose a smooth
nonnegative $K$ of integral one, supported in a sufficiently small
fixed neighborhood of zero that $|\widehat K(\xi)|\geq1/2$ for
$|\xi|\leq1$, using the Fourier convention with $e(-\xi u)$.
Plancherel on $\R$ gives
\begin{equation}\label{eq:local-convolution}
 \int_{|\beta|\leq1/H}|\widehat a(\beta)|^2\,d\beta
 \ll H^{-2}\int_{\R}
        \left|\sum_w a_w K((w-v)/H)\right|^2\,dv.
\end{equation}
Only $v\asymp Y$ contributes. Keep the integral restricted to this
range when replacing the kernel by $K(v\log(w/v)/H)$; no estimate
for this logarithmic kernel near $v=0$ is used. On the union of
the two kernels' supports within this range,
$|w-v|\ll H$ and their arguments differ by $O(H/Y)$. Therefore,
by Cauchy--Schwarz over $O(H)$ integers and then integrating the
centers allowed for each $w$, the squared norm of the error in
\Cref{eq:local-convolution} is
\begin{equation}\label{eq:kernel-error}
 \ll(H/Y)^2\sum_w|a_w|^2.
\end{equation}
This is smaller than $YL^{-A}$ for every fixed $A$, since
$H=x^{o(1)}$ and $Y=x^{1+o(1)}$.

Set $h=H/Y$ and $P(t)=\sum_w a_ww^{it}$. With $v=Ye^u$,
logarithmic Fourier inversion gives
\[
 \sum_w a_wK(v\log(w/v)/H)
  =\frac h{2\pi}\int_{\R}
         e^{-u}\widehat K(he^{-u}t/(2\pi))v^{-it}P(t)\,dt.
\]
Insert a smooth compactly supported function $\rho(u)$ equal to
one on the required range. The Fourier transform in $u$ of
$\rho(u)e^{-u}\widehat K(he^{-u}t/(2\pi))$, denoted $A(s,t)$,
satisfies, for each fixed $j$,
\[
 |A(s,t)|\ll_j(1+|s|)^{-2}(1+|ht|)^{-j}.
\]
This follows by two integrations by parts in $u$, since
$\widehat K$ is Schwartz and $u$ stays in a fixed compact interval.
Fourier-expand in $u$, apply Minkowski's integral inequality in $s$,
and apply Plancherel in $u$ for each $s$. As $dv\ll Y\,du$, the
normalization is $H^{-2}Yh^2=Y^{-1}$. After replacing $j$ by a
larger integer, we obtain
\begin{equation}\label{eq:local-mellin-energy}
 \int_{|\beta|\leq1/H}|\widehat a(\beta)|^2\,d\beta
 \ll_j Y^{-1}\int_{\R}(1+|ht|)^{-j}|P(t)|^2\,dt
          +(H/Y)^2\sum_w|a_w|^2.
\end{equation}

By \Cref{lem:dirichlet-mean}, for any dyadic time scale $T'\geq1$,
\[
 Y^{-2}\int_{T'\leq|t|\leq2T'}|P(t)|^2\,dt
          \ll (1+T'/Y)L^{C_9}.
\]
Consequently the portion $|t|>L/h$ in
\Cref{eq:local-mellin-energy} is $O(YL^{-A})$ after choosing $j$
large enough. For the remaining portion split the four factor
variables in \Cref{eq:four-factor-endpoint} into dyads. There are
$O(L^4)$ boxes and only boxes whose joint size is comparable to $Y$
contribute. Write
\[
 \frac{P(t)}Y=\sum_w\frac{1}{w}
       \left((w/Y)\Psi_1(w/Y)\right)w^{i(t+v_1)}
       \sum_{mp_sp r=w}(\text{the four coefficients}).
\]
Fourier-separate the parenthesized cutoff in logarithmic coordinates.
Its integral cost is a fixed log power; truncate at a fixed
log-power frequency with arbitrary saving. The tails are bounded
by the preceding mean values. The frequency shifts enlarge
$|t|\leq L/h$ only to $|t|\leq2L/h$, because $1/h$ exceeds every
fixed log power. We have reduced the assertion to
\begin{equation}\label{eq:polynomial-target}
 \int_{|t|\leq2L/h}|P_s(t)P_l(t)\mathcal R(t)|^2\,dt
                                      \ll L^{-A'}
\end{equation}
for any prescribed fixed $A'$, on each retained box, where
\begin{align*}
 P_s(t)&=\sum_{p_s\asymp P}b_s(p_s)p_s^{-1+it},\\
 P_l(t)&=\sum_{\substack{p\asymp P',\ p\in I_p\\\Pminus(p)>W}}
                       \chi_p(p)p^{-1+i\sigma_1+it},\\
 \mathcal R(t)&=
     \left(\sum_{m\asymp M_0}\alpha_m\chi_m(m)m^{-1+i\sigma_0+it}\right)
     \left(\sum_{r\asymp R_0}d_rr^{-1+it}\right).
\end{align*}
Here and below dyads may be intersected with their existing support.
We have
\begin{equation}\label{eq:four-factor-scales}
 e^{L^a}/2\leq P\leq e^{L^b},\qquad
 x^\tau/2\leq P'\leq x^\eta,\qquad
 PM_0P'R_0\asymp Y.
\end{equation}
For any subproduct of these four polynomials, collapse its coefficients
to $\sum_{n\asymp U}c_nn^{it}$. A fixed number of convolution factors
and \Cref{lem:divisor-moments} give
\begin{equation}\label{eq:collapsed-polynomial-l2}
       \sum_n|c_n|^2\ll L^{C_{10}}/U.
\end{equation}
This uses a fixed divisor moment: the number of marked groups has
already been absorbed by \Cref{eq:mark-sup}, not by a growing divisor
order. Individually all these polynomials are also $L^{O(1)}$ in
absolute value by their harmonic sums.

The remaining integral has three regimes. Most times are controlled
by the small-prime polynomial $P_s$: where $|P_s(t)|$ is a sufficiently
small negative power of $L$, the ordinary mean square of $P_l\mathcal R$
suffices. We will show that the exceptional times occupy only
$Y^{o(1)}$ unit intervals. On these intervals at large times,
$P_l$ is small by oscillation over its long rough-integer range;
a sparse mean square for $\mathcal R$ makes that saving sufficient.
For $|t|$ bounded by a fixed power of $L$, the $m$ factor in
$\mathcal R$ is small by the original discrepancy hypothesis.
We prove the exceptional-time and long-factor estimates first, then
choose the discrepancy precision in the completion of the argument.

\subsection{Exceptional times and a sparse mean square}

The split into ordinary and exceptional times, followed by high
prime-polynomial moments and a sparse mean-square estimate, is in the
spirit of the method of Matom{\"a}ki and Radziwi{\l}{\l};
compare \cite[Section~2.1 and Section~4, Lemmas~8--9]{MatomakiRadziwill2016}.
The estimates needed here are proved below; we do not invoke their
short-interval theorem.

Fix a large constant $A_s$ and let
$\mathcal E=\{t:|t|\leq2L/h,\ |P_s(t)|>L^{-A_s}\}$.
Outside $\mathcal E$, use the indicated gain and the mean square
of $P_l\mathcal R$, of joint size $Y/P$. Its length exceeds the
time range, since
\[
 \frac{Y/P}{2L/h}=\frac H{2LP}\longrightarrow\infty
\]
by $b<c$ and \Cref{eq:comparison-scales,eq:four-factor-scales}.
Thus this portion of \Cref{eq:polynomial-target} is
$O(L^{-2A_s+C_{11}})$.

The factorial coefficient estimate in the next proof is the standard
prime-polynomial moment argument; compare
\cite[Lemma~3 and its proof]{Soundararajan2009}. We keep its dependence
on the growing power explicit.

\begin{lemma}[Number of exceptional unit intervals]\label{lem:exceptional-times}
The set $\mathcal E$ meets at most
\[
 \exp\!\left(O(L^b+L^{1-a}\log L)\right)=Y^{o(1)}
\]
intervals $[j,j+1]$, uniformly in the bounded coefficients $b_s$.
\end{lemma}

\begin{proof}
Choose a point of $\mathcal E$ in each occupied interval, and split
the interval indices into three residue classes to obtain separated
points. Put $j_0=\lfloor\log Y/\log(2P)\rfloor$. The polynomial
$P_s(t)^{j_0}$ has indices at most $Y$. Its coefficient-square sum
is at most
\[
 j_0!\left(\sum_{p_s\asymp P}|b_s(p_s)|^2p_s^{-2}\right)^{j_0}
                  \leq j_0!(C/P)^{j_0}.
\]
Indeed an integer has at most $j_0!$ ordered representations as
a product of $j_0$ primes; Cauchy--Schwarz on each fiber proves
the first inequality even with repeated primes.

For a polynomial $Q(t)=\sum_{n\leq Y}c_nn^{it}$ and unit-separated
points in $[-O(Y),O(Y)]$, the unit-interval Sobolev inequality and
bounded overlap imply
\[
 \sum_j|Q(t_j)|^2\ll\int_{-O(Y)}^{O(Y)}(|Q(t)|^2+|Q'(t)|^2)\,dt
                  \ll Y(\log(2Y))^{O(1)}\sum_n|c_n|^2.
\]
In the last step use \Cref{lem:dirichlet-mean}; differentiation
multiplies a coefficient by $i\log n$, so the constants here are
independent of $j_0$. The range $2L/h$ is $o(Y)$, as required.
If $N_{\mathcal E}$ is the number of occupied intervals, it follows
that
\[
 N_{\mathcal E}\ll YL^{O(1)}j_0!
                   (CL^{2A_s}/P)^{j_0}.
\]
Now $j_0=O(L^{1-a})$ and
$\log Y-j_0\log P=O(L^b+j_0)$ by the definition with $2P$.
Taking logarithms and using $\log(j_0!)\leq j_0\log j_0$ proves
the asserted estimate.
\end{proof}

\begin{lemma}[Sparse residual mean square]\label{lem:sparse-residual}
For the occupied unit intervals $I_j$ in
\Cref{lem:exceptional-times},
\[
             \sum_j\sup_{t\in I_j}|\mathcal R(t)|^2\ll L^{C_{12}}.
\]
\end{lemma}

\begin{proof}
The collapsed residual size is
\[
 U\asymp Y/(PP')\geq Y^{1-\eta-o(1)}>Y^{0.6}
\]
for large $x$, and its coefficient-square sum is $L^{O(1)}/U$ by
\Cref{eq:collapsed-polynomial-l2}. On a containing integer interval
$n\asymp U$, the Gram kernel satisfies, uniformly for $|z|\ll Y$,
\begin{equation}\label{eq:sparse-gram-kernel}
 \left|\sum_{n\asymp U}n^{iz}\right|
             \ll \frac U{1+|z|}+Y^{1/2}.
\end{equation}
For $|z|<1$ this is trivial. For $1\leq|z|\leq cU$ with small
fixed $c$, the phase $z\log n/(2\pi)$ has monotone first derivative
of magnitude comparable to $|z|/U$ and less than $1/2$;
\Cref{lem:derivative-tests} gives $O(U/|z|)$. For $cU<|z|\ll Y$
the second derivative test gives
$O(|z|^{1/2}+U|z|^{-1/2})=O(Y^{1/2})$.

Choose a maximum point of $|\mathcal R|$ in each closed $I_j$ and
again split the indices into three residue classes. In each class
these points are at least unit separated. For $R=Y^{o(1)}$ such
points, the absolute row sums of their Gram matrix are at most
\[
 C\bigl(U\log(2Y)+RY^{1/2}\bigr)\ll U\log(2Y).
\]
The first term follows by grouping distances into unit intervals;
the second is absorbed by $U>Y^{0.6}$. Schur's bound on this Gram
matrix, applied to the map
$(c_n)\mapsto(\sum_nc_nn^{it_j})_j$, now gives
\[
 \sum_j|\mathcal R(t_j)|^2
       \ll U\log(2Y)\sum_n|c_n|^2\ll L^{O(1)}.
\]
Summing the three classes proves the statement, including the
suprema over the intervals.
\end{proof}

\subsection{Cancellation of the long rough-integer factor}

The following elementary estimate is the reason for requiring a
rough integer factor of positive power length in every endpoint.

\begin{lemma}[Long rough-integer polynomial]\label{lem:long-rough-polynomial}
Fix $\tau>0$, $\eta<1/4$, $C>0$, and $A>0$. Let
$x^\tau/2\leq P'\leq x^\eta$, let $J'\subset[P',2P']$ be an interval,
and let $\chi$ have modulus at most $L^C$. If $|\sigma|\leq L^C$,
then, for a sufficiently large fixed $B_0$, uniformly whenever
$L^{B_0}\leq|t|$ and $|t+\sigma|\leq x^2$,
\[
 \left|\sum_{\substack{n\in J'\\\Pminus(n)>W}}
                \chi(n)n^{-1+i(t+\sigma)}\right|\ll L^{-A}.
\]
The constant $B_0$ can be chosen after $C,A,\tau,\eta$.
\end{lemma}

\begin{proof}
Set $h_s=2\lceil T^2\rceil$ and
$D=W^{4h_s+2}=\exp(O(\sqrt L\,T^2))=x^{o(1)}$.
By \Cref{lem:rough-sieve-polynomial}, there is a polynomial
\[
 U(n)=\sum_{d\mid n}\lambda_d\geq\ind_{\Pminus(n)>W},
\]
whose coefficients have modulus at most one, are supported on
squarefree $d\leq D$ with primes at most $W$, and satisfy
\[
 \sum_{n\in[P',2P']}\bigl(U(n)-\ind_{\Pminus(n)>W}\bigr)
           \ll P'e^{-h_s}+D.
\]
The summand is nonnegative, so restriction to any $J'$ preserves
this upper bound. Replacing roughness by $U$ in the normalized
sum therefore costs at most $O(e^{-h_s}+D/P')$, smaller than every
fixed negative power of $L$. Also
\begin{equation}\label{eq:sieve-harmonic-cost}
 \sum_d\frac{|\lambda_d|}{d}
       \leq\prod_{p\leq W}(1+1/p)\ll\log W=\sqrt L.
\end{equation}

Write $q'$ for the character modulus. Terms with $(d,q')>1$ vanish.
For the other $d$, split $n=dv$ by $v=q'y+a$ modulo $q'$. The
character is constant on a class. Its $y$ variable ranges over an
interval at size
$N=P'/(dq')\geq x^{\tau/2}$ for large $x$.
Put $u=t+\sigma$. By taking $B_0>C+1$ we ensure
$|u|\geq L^{B_0}/2$. We claim, on every subinterval of the allowed
$y$ range,
\begin{equation}\label{eq:log-phase-sum}
 \left|\sum_y\exp(iu\log(q'y+a))\right|
                 \ll N\bigl(|u|^{-1}+N^{-\delta}\bigr)
\end{equation}
for a fixed $\delta>0$ depending only on $\tau$.

If $|u|\leq cN$, the monotone first derivative test gives
$O(N/|u|)$, since the derivative of $u\log(q'y+a)/(2\pi)$
stays away from nonzero integers. Otherwise put $\eps_0=1/10$.
The overlapping ranges
\[
       N^{k-2+\eps_0}\leq|u|\leq N^{k-\eps_0},\qquad k\geq2,
\]
cover $cN<|u|\leq x^2$ using
$2\leq k\leq K=\lceil4/\tau\rceil+3$. The $k$th derivative of
this phase has constant sign and magnitude comparable to
$|u|N^{-k}$. For $k=2$, the second derivative test directly saves
a positive power. For $k>2$, apply the van der Corput differencing
inequality $k-2$ times with positive integer shifts at most
$H_0=\lfloor N^{\eps_0/(2k)}\rfloor$.
After shifts $h_1,\ldots,h_{k-2}$, the second derivative is an
iterated integral of the $k$th derivative, and hence has constant
sign and magnitude comparable to
\[
 |u|N^{-k}h_1\cdots h_{k-2},
 \quad\text{between }c_kN^{-2+\eps_0}
                      \text{ and }C_kN^{-\eps_0/2}.
\]
On an interval of any length $m\leq C N$, the second derivative
bound is
\[
 O\left(m\sqrt\lambda+\lambda^{-1/2}\right)
                         \ll N^{1-\eps_0/4}.
\]
To justify the estimate uniformly for short subintervals, extend
each sequence by zero in a containing interval of length $O(N)$.
The differencing inequality for its normalized sum has the form
$B_j^2\ll H_0^{-1}+B_{j+1}$, where $B_{j+1}$ bounds the average
of the absolute normalized correlations. Their supports are
intersections of translates of the original interval, and so are
again intervals. All shifts are $o(N)$, preserving the derivative
bounds. Iterating from the final exponent $\eps_0/4$ proves a
bound $O(N^{-\delta_k})$ for the original normalized sum, with
$\delta_k=\eps_0/(k2^{k-1})$. Taking
$\delta=\eps_0/(K2^{K-1})$ proves \Cref{eq:log-phase-sum} in every
case. The differencing inequality itself follows by averaging
$H_0$ translates and applying Cauchy--Schwarz; its zero-shift term
is precisely the displayed $O(H_0^{-1})$.

Partial summation of $1/[d(q'y+a)]$ costs $O(1/P')$ times the
uniform unweighted bound. Summing the at most $q'$ classes gives
\[
 \left|\sum_{\substack{n\in J'\\d\mid n}}
                      \chi(n)n^{-1+iu}\right|
   \ll\frac1d\bigl(|u|^{-1}+x^{-\tau\delta/2}\bigr).
\]
Finally \Cref{eq:sieve-harmonic-cost} bounds the total by
\[
 O(e^{-h_s}+D/P')+
 O\left(\sqrt L\,(L^{-B_0}+x^{-\tau\delta/2})\right).
\]
Choosing $B_0>A+2$ in addition to its previous constraint proves
the lemma.
\end{proof}

\subsection{Completion and order of parameters}

We finish \Cref{eq:polynomial-target}. Outside $\mathcal E$ its
integral is $O(L^{-2A_s+C_{11}})$, so choose $A_s$ sufficiently
large. On the exceptional intervals with $|t|\geq L^{B_0}$,
\Cref{lem:long-rough-polynomial} gives an arbitrarily strong
uniform bound for $P_l$. Its range condition holds since
$2L/h=o(Y)=x^{1+o(1)}<x^2/2$; the fixed frequency shift does not
change this. The polynomial $P_s$ is bounded, and
\Cref{lem:sparse-residual} gives
\[
 \int_{\mathcal E\cap\{|t|\geq L^{B_0}\}}
           |P_s(t)P_l(t)\mathcal R(t)|^2\,dt
 \ll \bigl(\sup_{|t|\geq L^{B_0}}|P_l(t)|^2\bigr)L^{O(1)}.
\]
Here the supremum is restricted to the actual time range, and
each occupied interval has length one. Choose the saving in the
long-polynomial lemma after the exponent in the sparse bound.

For all remaining $|t|\leq L^{B_0}$ use
\Cref{eq:discrepancy} on the $m$ polynomial in $\mathcal R$.
The retained dyads satisfy $PM_0P'R_0\asymp Y$, so
$m\ll Y=x^{1+o(1)}<x^2$ for large $x$, also after the smooth
cutoff has been separated. Its character modulus is a fixed
log power, and its frequency is $t+\sigma_0$. A single sufficiently
large $B$ therefore bounds this polynomial by $L^{-B}$ throughout
the range. Every other factor has absolute value $L^{O(1)}$, and
the length of this time interval is $2L^{B_0}$. The resulting
integral is $L^{-2B+B_0+O(1)}$, giving the required accuracy.
This proves \Cref{eq:polynomial-target}, then
\Cref{eq:local-energy-target} by \Cref{eq:local-mellin-energy},
and then every comparison bound. The residual pairing and the
raw identity in \Cref{eq:stripped-pattern} prove
\Cref{eq:shift-bound}.

For completeness, the choices occur in the following order.
First fix the endpoint bounds, group geometry, cutoff derivative
bounds, and desired saving $D_0$. The uniform endpoint norm bound
fixes the physical transfer target and its required ideal target;
the fixed Mellin cutoff fixes the ideal frequency range. Next choose
the probe count and comparison parameters of \Cref{thm:ideal},
then choose the fixed $J$ required by transference. All pattern,
cell, and absolute norm costs are now determined. Choose the
separation accuracies, arc exponent, local energy target, and
Mellin tail precision. Next choose $A_s$, the accuracy of the
long rough polynomial and its $B_0$. Finally choose one $B$ in
\Cref{eq:discrepancy} exceeding all the modulus, frequency, and
low-time saving requirements, and let $x$ tend to infinity.
The residual $L^2$ bounds used before the choice of $J$ depend only
on the endpoint coefficient bounds and a fixed divisor moment.
Increasing later Fourier tail accuracies uses more integrations
by parts, with no enlargement of the already fixed ideal frequency
range. Thus none of these choices is circular.
